


Generated by \texttt{python -m crashguard.rd\_docs\_b}. Preregistration: \texttt{docs/02\_preregistration/PREREGISTRATION\_V12.md} (addenda A and B, with the disclosures of section 10.6). Constants: \texttt{docs/03\_results/ASSUMPTIONS.md}. Session 1 results: \texttt{REDESIGN\_A\_RESULTS.md}. Code checks: \texttt{CODE\_VERIFICATION.md}.

\subsection*{1. Final Holm table (one family of ten)}

{\scriptsize\raggedright
\begin{longtable}{p{0.237\linewidth}p{0.237\linewidth}p{0.237\linewidth}p{0.237\linewidth}}
\toprule
Id & Raw p & Final Holm p & Verdict \\
\midrule
\endhead
H-D1a & 0.0005 & 0.004 & SUPPORTED \\
H-D1b & 0.0005 & 0.004 & SUPPORTED \\
H-D1c & 0 & 0 & SUPPORTED \\
H-D1s & 0.0005 & 0.004 & SUPPORTED \\
H-D2a & 0.0005 & 0.004 & SUPPORTED \\
H-D2b & 1 & 1 & NOT SUPPORTED \\
H-D2c & 0 & 0 & SUPPORTED \\
H-D3 & 1 & 1 & NOT SUPPORTED \\
H-D4a & 1 & 1 & NOT SUPPORTED \\
H-D4b & 1 & 1 & NOT\_RUN \\
\bottomrule
\end{longtable}}

Supported: H-D1a, H-D1b, H-D1c, H-D1s, H-D2a, H-D2c. Not supported or not run: H-D2b, H-D3, H-D4a, H-D4b. H-D1c and H-D2c are formula or bounded-model checks (p = 0 when met).

\subsection*{2. D3 (gated slotted retries), FRESH-B, next to the original GS-CA failure}

Original failure: GS-CA lost in all 16 burst cells of the old held-out grid; 0.54 of its failed Critical alerts failed through a collision or a deaf gateway (\texttt{DIAGNOSIS\_GSCA.md}). D3 (V12 10.2) slots only the retries, with the burst size taken from overheard frames, on top of the frozen D1 + D2. Grid: 32 burst cells drawn from the freeze hash, 300 hours per cell, run once.

{\scriptsize\raggedright
\begin{longtable}{p{0.475\linewidth}p{0.475\linewidth}}
\toprule
Quantity & Value \\
\midrule
\endhead
Critical delivered within 10 s, D1 + D2 (mean over cells) & 0.270 \\
same, D3 & 0.270 \\
D3 minus D1 + D2 (pooled, 95\% interval; 5 points required) & +0.001 [+0.000, +0.001] \\
cells with a gain of at least 5 points & 0 of 32 \\
smallest per-cell difference in delivery before the own deadline (no-harm; lower bound) & -0.001 (-0.004) \\
runs above the duty limit & 0 \\
\bottomrule
\end{longtable}}

\textbf{D3 is a negative result} (kill rule of V12 10.2: the 5-point gain is not met; the no-harm part holds). The gain is real in sign (interval above zero) and about 80 times smaller than required. \textbf{Mechanism seen on the DEV grid (not on FRESH-B):} in a DEV burst hour the slots were used 2.1 times in total for 25 burst alerts with up to 12 retries each, because a vehicle rarely decodes two other vehicles in this channel model (vehicle-to-vehicle links are mostly out of range), so a burst-size estimate without an oracle is rarely available. GS-CA and the Tamang-style baseline both used a burst size given by an oracle. This reading is inferred from counts, not isolated by an ablation. Design deviations that may matter are in V12 10.6 (the gate counts other vehicles; slots apply to the flash of a round; the 5 s epoch cap does not hold at SF12).

\subsection*{3. D4 (24-byte provisional frame), H-D4a, next to the original R2 and R3 results}

Original: R2 met the 5-point threshold in 10 of 16 cells but its sign test had p = 0.23 (power 0.22); R3's extra airtime was 38\% against a 40\% bound. The task text assumed R3's provisional frame was 97 bytes; in this repository it is a 14-byte flash (V12 10.1), so a 24-byte frame is longer, not shorter. H-D4a is the extra airtime of the two-stage scheme over the single-stage alert on the stored test events at the 1\% target, 48 cells (6 spreading factors x 8 loss levels drawn after the freeze); criterion: upper 95\% bound below 0.25 in every cell.

{\scriptsize\raggedright
\begin{longtable}{p{0.237\linewidth}p{0.237\linewidth}p{0.237\linewidth}p{0.237\linewidth}}
\toprule
Provisional frame & Smallest extra airtime over the 48 cells & Largest & Cells with upper bound below 0.25 \\
\midrule
\endhead
14-byte flash (R3 as registered) & 0.315 & 0.415 & 0 of 48 \\
24-byte frame, 8-byte tag (D4) & 0.410 & 0.496 & 0 of 48 \\
97-byte frame (task wording of R3) & 1.076 & 1.140 & 0 of 48 \\
\bottomrule
\end{longtable}}

\textbf{D4 is a negative result.} H-D4a fails in every cell, and the 24-byte frame costs more airtime than R3's own flash, as the arithmetic of V12 10.1 said. H-D4b (arrival time and burst delivery) was not run: it needs a two-stage alert with per-alert frame lengths in the simulator, which does not exist, and the kill rule ends D4 when H-D4a fails. NOT\_RUN; it enters the family at p = 1. The lesson for the paper is about R3's premise: its provisional frame was already short.

\subsection*{4. Firmware decision path (D6), host build}

C decision path (\texttt{firmware\_rd/}): time on air, SHA-256 and HMAC, the 97-byte alert and its tag, freshness, the D1 table lookup and Beta counters, Ed25519 verification of the cluster ACK (TweetNaCl, public domain), the D2 decision and defer rule, and the D3 slot rule. Differential test against the Python reference used by the simulator: 10000 random vehicle histories, 409999 compared operations (numbers bit for bit, frames and tags byte for byte), \textbf{0 mismatches}. Replay of 5140 decisions recorded inside real simulator runs (D1 lookups, defers, D3 slots, D2 decisions): \textbf{0 mismatches}. Host-build object sizes (x86-64, -Os, not the ESP32): firmware logic 5985 bytes of code, whole TweetNaCl 12608 bytes; D1 table 186624 bytes per base spreading factor. \textbf{ESP32 flash and RAM, and the emulator run: NOT\_RUN} (no Xtensa or RISC-V toolchain, no Espressif QEMU, no ESP-IDF here). The statement allowed in the paper: "the firmware decision path was executed on simulated hardware-format inputs in a host build; an ESP32 emulator and a physical device were not used; no measurement on physical hardware is reported." The word "simulated hardware-format inputs" refers to the recorded simulator decisions and random traces; the hardware-format logs of \texttt{data/simulated\_hardware} were not replayed through the C path (NOT\_RUN). The trigger and severity code are the earlier exports (\texttt{firmware\_export/}) and were not re-run in D6. Double precision is used so that results match Python bit for bit; an ESP32 port needs float or fixed point (NOT\_RUN).

\subsection*{5. Cross-simulator check for D1 (X1.3)}

Full tables and the list of differences: \texttt{XSIM\_VERDICTS.md}. Summary on scenario S-X (32 SF9 cells of FRESH-A, 100 hours, ns-3 lorawan module and the custom simulator, frozen D1):

{\scriptsize\raggedright
\begin{longtable}{p{0.237\linewidth}p{0.237\linewidth}p{0.237\linewidth}p{0.237\linewidth}}
\toprule
Hypothesis & Custom simulator & ns-3 & Same \\
\midrule
\endhead
H-D1a & met & not met & no \\
H-D1b & met & not met & no \\
H-D1c & met & met & yes \\
H-D1s & not met & not met & yes \\
\bottomrule
\end{longtable}}

D1 minus the severity table at loss of 40\% or more: custom +0.116 [+0.113, +0.120], ns-3 +0.021 [+0.017, +0.025]. \textbf{D1's gain over the severity table is simulator-dependent}: it holds in the custom simulator and does not reach the 5-point bar or the no-harm bar in ns-3. It cannot be in the abstract, and D1 does not enter the main paper (V12 section 8 and D8). The calibration criterion of X1.2 was not met (one cell 5.4 points out), so even the agreement on H-D1c is stated under that condition. Against the best baseline of each cell (H-D1s) D1 is not supported in either simulator on S-X (on FRESH-A it was supported; the baselines there were chosen per cell on the same data). Possible reasons not isolated: the custom policies use a flash frame whose reception counts as delivery and ns-3's do not; different ACK length, timing and collision models.

\subsection*{5b. Added after the second Session 2 prompt (ADDED-AFTER-SESSION-1)}

Nessa-style baseline, calibration under two sensitivity settings and the ns-3 run with the new baseline: \texttt{XSIM\_VERDICTS.md}, \texttt{XSIM\_CALIBRATION.md}, \texttt{REDESIGN\_DECISION\_TABLE.md}; keys \texttt{rd.fc.*}, \texttt{rd.x1.b.*}, \texttt{rd.x13.*}. None of this is in the Holm family and none changes a verdict above.

\subsection*{6. Not done in session 2 (NOT\_RUN, with reasons)}

\begin{itemize}
\item D4b (needs a feature the simulator lacks; moot after the failed H-D4a); ESP32 size and emulator (no toolchain); hardware-format log replay through the C path; ADR-on baseline in ns-3 (no data for ADR in a one-hour run); X1.4 comparison with the three published simulator papers (none opened); implementation of an adaptive scheme from the five closest papers as an extra baseline (none could be implemented from the text read: \texttt{CLOSEST\_PRIOR\_WORK.md}); the bench run of the two-stage alert (the author's: \texttt{docs/08\_hardware/R3\_BENCH\_PROTOCOL.md}; \texttt{rd.x4.hw\_row} NOT\_RUN).
\item Searches: OpenAlex returned HTTP 503; IEEE Xplore, Scopus, Google Scholar and TRID were not run (\texttt{NOVELTY\_SCAN.md}).
\end{itemize}

